% Generated draft: requires mathematical and source review.
\documentclass[11pt]{article}
\usepackage[margin=25mm]{geometry}
\usepackage{fontspec}
\setmainfont{Latin Modern Roman}
\usepackage{amsmath,amssymb,mathtools,unicode-math}
\setmathfont{Latin Modern Math}
\usepackage{enumitem,needspace,xurl,hyperref}
\setlength{\parindent}{0pt}
\setlength{\parskip}{5pt}
\setlength{\emergencystretch}{3em}
\newcommand{\N}{\mathbb{N}}
\newcommand{\Z}{\mathbb{Z}}
\newcommand{\Q}{\mathbb{Q}}
\newcommand{\R}{\mathbb{R}}
\newcommand{\C}{\mathbb{C}}
\newcommand{\F}{\mathbb{F}}
\newcommand{\A}{\mathbb{A}}
\newcommand{\PP}{\mathbb{P}}
\newcommand{\E}{\mathbb{E}}
\newcommand{\eps}{\varepsilon}
\newcommand{\dd}{\,\mathrm{d}}
\begin{document}
\section*{UnsolvedMath ID: 800646}
\textbf{Source problem code:} catalogue entry
\subsection{Definitions and mathematical statement}
For a set \(X\), a cover of \(X\) is a family of subsets whose union is \(X\). A finite subcover is a finite subfamily whose union is \(X\). For specified families \(\mathcal A\) and \(\mathcal B\) of covers of \(X\), write
\[
\binom{\mathcal A}{\mathcal B}
\quad\Longleftrightarrow\quad
(\forall\mathcal U\in\mathcal A)(\exists\mathcal V\subseteq\mathcal U)\;\mathcal V\in\mathcal B.
\]
The selection principle \(\mathsf U_{\mathrm{fin}}(\mathcal A,\mathcal B)\) means that for every sequence \((\mathcal U_n)_{n\in\N}\) of members of \(\mathcal A\), each having no finite subcover, there exist finite subfamilies \(\mathcal F_n\subseteq\mathcal U_n\) such that
\[
\left\{\bigcup\mathcal F_n:n\in\N\right\}\in\mathcal B.
\]
Here \(\N=\{1,2,3,\ldots\}\).

The supplied source record asks whether the following two equivalences hold, with all cover families interpreted on the same underlying space \(X\):
\[
\mathsf U_{\mathrm{fin}}\bigl(\Gamma(X),\Omega(X)\bigr)
\quad\Longleftrightarrow\quad
\binom{\Lambda(X)}{\mathcal O^{\omega\text{-}*\mathrm{gp}*}(X)},
\]
\[
\mathsf U_{\mathrm{fin}}\bigl(\Gamma(X),\mathrm T(X)\bigr)
\quad\Longleftrightarrow\quad
\binom{\Lambda(X)}{\mathcal O^{\tau\text{-}*\mathrm{gp}*}(X)}.
\]
These are two separate questions attributed by the registry to Tsaban, 2003, Problem 19.

\emph{Unresolved qualifications.} The supplied record does not specify the class of spaces over which either equivalence is to hold. It also does not define \(\Gamma\), \(\Omega\), \(\Lambda\), \(\mathrm T\), or the starred superscripts in \(\mathcal O^{\omega\text{-}*\mathrm{gp}*}\) and \(\mathcal O^{\tau\text{-}*\mathrm{gp}*}\). In particular, the meanings of the starred groupability notation, restrictions on cover cardinality, openness requirements, and conventions concerning a cover containing \(X\) must be recovered from the source before this can be a fully specified mathematical statement. The definitions of the two operators above give their usual logical forms; they do not supply these missing cover classes or domain hypotheses. No interpretation of the stars is assumed.
\subsection{Short English statement}
Tsaban's Problem 19 asks two equivalence questions: whether \(\mathsf U_{\mathrm{fin}}(\Gamma,\Omega)\) and \(\mathsf U_{\mathrm{fin}}(\Gamma,\mathrm T)\) respectively coincide with the two displayed subcover properties involving starred groupability notation. The intended spaces and cover definitions are missing from the supplied record, so the formulation remains conditional on recovering them. The registry labels the problem open; this is supplied status information, not an independently verified assessment.
\subsection{Sources}
\begin{itemize}
\item Supplied local UnsolvedMath registry record: permanent ID \(800646\), exact source problem code \texttt{catalogue entry}, title ``Selection Principles Milestone Problems --- Problem 19,'' in the AMR Open Problem Lists collection. The record attributes the problem to Tsaban, with proposed year 2003, and labels its status ``open.''
\item Supplied source reference: Tsaban, \emph{Selection Principles in Mathematics: A Milestone of Open Problems} (2003), Problem 19. Supplied URL: \url{https://arxiv.org/abs/math/0312182}.
\item Supplied provenance reports access on 2026-10-16 and extraction method \texttt{arxiv\_latex}. These are source-record provenance claims; the cited source and its definitions have not been independently inspected for this draft.
\end{itemize}
\end{document}
